% Qiushi Engine default English research-paper template.
% Target: APS Physical Review Letters in REVTeX 4.2 reprint (two-column) form.
% Build with: latexmk -lualatex -interaction=nonstopmode -halt-on-error main.tex
% If the LuaLaTeX font cache is not writable, prefix that command with
% TEXMFVAR=/tmp/qiushi_texmf_var (or another session-writable cache directory).
%
% Keep this class line, page geometry, typography, and column behavior intact for
% PRL output. Add field-specific packages only when the manuscript genuinely needs
% them and compile immediately after doing so. Keep figures under figures/ or
% ../../workspace/figures/, and keep verified bibliography records in refs.bib.
%
% Before final delivery: remove every placeholder; resolve all citations and cross
% references; inspect the log for errors, missing assets/glyphs, stuck floats, and
% overfull content; then render and inspect every page at its final size. A clean
% compile alone does not establish readable figures, tables, equations, or pages.
\documentclass[aps,prl,reprint,superscriptaddress,nofootinbib,longbibliography]{revtex4-2}

\usepackage{graphicx}
\usepackage{amsmath,amssymb,mathtools,bm}
\usepackage{amsthm}
\usepackage{siunitx}
\usepackage{booktabs}
\usepackage{tabularx}
\usepackage{array}
\usepackage{xurl}
\usepackage{xcolor}
\usepackage{microtype}
\usepackage[colorlinks=true,linkcolor=blue!55!black,citecolor=blue!55!black,urlcolor=blue!55!black]{hyperref}

\graphicspath{{figures/}{../figures/}{./}{../../workspace/figures/}}
\sisetup{detect-all=true,per-mode=symbol}
\Urlmuskip=0mu plus 2mu
\emergencystretch=1.5em

\providecommand{\dd}{\mathop{}\!\mathrm{d}}
\providecommand{\dv}[2]{\frac{\dd #1}{\dd #2}}
\providecommand{\pdv}[2]{\frac{\partial #1}{\partial #2}}
\providecommand{\vect}[1]{\boldsymbol{#1}}
\providecommand{\mat}[1]{\mathbf{#1}}
\providecommand{\E}{\mathbb{E}}
\providecommand{\Var}{\operatorname{Var}}
\providecommand{\argmin}{\operatorname*{arg\,min}}
\providecommand{\abs}[1]{\lvert#1\rvert}
\providecommand{\norm}[1]{\lVert#1\rVert}

\newtheorem{theorem}{Theorem}
\newtheorem{proposition}{Proposition}
\newtheorem{lemma}{Lemma}
\newtheorem{corollary}{Corollary}
\theoremstyle{definition}
\newtheorem{definition}{Definition}
\theoremstyle{remark}
\newtheorem{remark}{Remark}

% Layout contract:
% - Use figure/table when the complete object remains legible at one-column width.
% - Use figure*/table* for genuinely dense multi-panel or comparison objects that
%   need both columns. Do not make every object full width by habit.
% - Prefer vector PDF for plots and line art; use sufficient-resolution raster media
%   when the source is intrinsically raster. All labels must remain readable at the
%   final printed size, with consistent symbols, units, panel order, and captions.
% - Build multi-panel graphics as one intentional scientific composition. Avoid
%   subfigure packages that fight REVTeX; place panel labels in the source artwork.
% - Break equations semantically with aligned, split, or multline first. Use
%   widetext only when one inseparable mathematical structure needs full page width.
% - Use ruledtabular/booktabs conventions, explicit units and uncertainties, and no
%   decorative vertical rules. Redesign a crowded table instead of shrinking it.
% - Never conceal crowding with resizebox, arbitrary negative spacing, forced [H]
%   placement, altered margins, or unreadably small text.
% - Section structure follows the scientific argument. Rename, merge, or omit the
%   sample sections as the work requires; do not preserve headings as empty shells.

\begin{document}

\title{Article Title}

\author{Qiushi Engine}
\noaffiliation
% Add \affiliation{...}, \email{...}, or additional authors only when the
% corresponding real, user-approved metadata is available. Never repeat the
% author name as a fictional affiliation.

\date{\today}

\begin{abstract}
Abstract text.
\end{abstract}

\maketitle

\section{Introduction}

Text.

\section{Methods}

Text.

\begin{equation}
  \begin{aligned}
    \mathcal{L}u &= f, \\
    \mathcal{M}v &= g .
  \end{aligned}
  \label{eq:column-example}
\end{equation}

\begin{figure}[t]
  \centering
  % Replace the box with \includegraphics[width=\columnwidth]{figure_1.pdf}.
  \fbox{\rule{0pt}{1.25in}\rule{0.88\columnwidth}{0pt}}
  \caption{State what is shown, define encodings and conditions, and explain the
  result needed for the argument.\label{fig:single-column}}
\end{figure}

\section{Results}

Text.

\begin{figure*}[t]
  \centering
  % Replace the box with \includegraphics[width=\textwidth]{figure_2.pdf}.
  \fbox{\rule{0pt}{1.55in}\rule{0.96\textwidth}{0pt}}
  \caption{Use both columns only when the full scientific composition and its
  labels cannot remain readable in one column.\label{fig:two-column}}
\end{figure*}

\begin{widetext}
\begin{equation}
  \mathcal{K}(x,y;t)
  = \int_{\Omega} A(x,z;t)B(z,y;t)\,\dd z
  + \int_{0}^{t}\!\int_{\Omega} C(x,z;t-s)D(z,y;s)\,\dd z\,\dd s .
  \label{eq:full-width-example}
\end{equation}
\end{widetext}

\begin{table}[b]
  \caption{A compact one-column comparison with quantities, units, and uncertainty
  defined in the caption or headings.\label{tab:single-column}}
  \begin{ruledtabular}
  \begin{tabular}{lcc}
  Case & Value A & Value B \\
  \colrule
  A & -- & -- \\
  B & -- & -- \\
  \end{tabular}
  \end{ruledtabular}
\end{table}

\begin{table*}[t]
  \caption{A comparison whose natural width and information density genuinely
  require both columns.\label{tab:two-column}}
  \begin{ruledtabular}
  \begin{tabular*}{\textwidth}{@{\extracolsep{\fill}}lccccc@{}}
  Method & Assumption & Observable & Control & Uncertainty & Result \\
  \colrule
  A & -- & -- & -- & -- & -- \\
  B & -- & -- & -- & -- & -- \\
  \end{tabular*}
  \end{ruledtabular}
\end{table*}

\section{Discussion}

Text.

\section{Conclusion}

Text.

\begin{acknowledgments}
Acknowledgment text.
\end{acknowledgments}

\IfFileExists{refs.bib}{\bibliographystyle{apsrev4-2}\bibliography{refs}}{}

\end{document}
